\hypertarget{ej2}{}
\noindent \textbf{Ejercicio 2.} Considere la siguiente oración: ``Si es mi cumpleaños o hay torta, entonces hay torta''.
\begin{itemize}
    \item Escribir usando lógica proposicional y realizar la tabla de verdad
    \item Asumiendo que la oración es verdadera y hay una torta, ¿qué se puede concluir?
    \item Asumiendo que la oración es verdadera y no hay una torta, ¿qué se puede concluir?
    \item Suponiendo que la oración es mentira (es falsa), ¿se puede concluir algo?
\end{itemize}

\vspace{0.2cm}
{\color{gray}\hrule}
\vspace{0.4cm}

\textbf{Resolución:}

Definimos las proposiciones simples: \\
$a = \text{Es mi cumpleaños}$ \\
$b = \text{Hay torta}$

\begin{enumerate}[label=\textbf{\arabic*)}]
    \item \textbf{Escribir usando lógica proposicional y realizar la tabla de verdad} \\
    La oración se traduce como: $(a \vee b) \Rightarrow b$

    \begin{center}
    \begin{tabular}{|c|c|c|c|}
    \hline
    $a$ & $b$ & $(a \vee b)$ & $(a \vee b) \Rightarrow b$ \\
    \hline
    V & V & V & \textbf{V} \\
    V & F & V & \textbf{F} \\
    F & V & V & \textbf{V} \\
    F & F & F & \textbf{V} \\
    \hline
    \end{tabular}
    \end{center}

    \item \textbf{Oración verdadera y hay torta} \\
    $(a \vee b) \Rightarrow b = \text{True}$, $b = V$. Entonces la oración es True, $b$ es True y $a$ no tiene un valor concreto, \underline{puede ser True o False}.

    \item \textbf{Oración verdadera y no hay torta} \\
    $(a \vee b) \Rightarrow b = \text{True}$, $b = F$. Entonces necesariamente el antecedente $(a \vee b) = F$, lo que implica que \underline{$a = \text{False}$}.

    \item \textbf{Oración es mentira (falsa)} \\
    $(a \vee b) \Rightarrow b = \text{False}$. Para que una implicación sea falsa, necesariamente el consecuente debe ser falso ($b = F$) y el antecedente verdadero. Entonces $(a \vee b) = \text{True}$, por lo tanto \underline{$a = \text{True}$}.
\end{enumerate}

\vspace{0.5cm}
\hrule
\vspace{0.5cm}